%%
%% Nature Machine Intelligence Manuscript Shell (Option B Fallback)
%% Venue: Nature Machine Intelligence (Springer Nature)
%% Submission Mode: Rolling submission (contingent on post-rejection activation from WWW'27)
%%
\documentclass[11pt,a4paper]{article}

\usepackage[utf8]{inputenc}
\usepackage{amsmath,amssymb,amsfonts}
\usepackage{graphicx}
\usepackage{booktabs}
\usepackage{lineno}
\usepackage{cite}
\usepackage{hyperref}
\usepackage[margin=2.5cm]{geometry}

\linenumbers

\title{\textbf{The download-weighted safety gap on public machine learning model registries is falsified on registered estimators}}

\author{
  Petar Radanliev$^{1,2,*}$\\[1ex]
  \small $^{1}$School of Computing and Communications, STEM, The Open University, Milton Keynes, UK\\
  \small $^{2}$Department of Computer Science, University of Oxford, Oxford, UK\\
  \small $^*$Corresponding author: \texttt{petar.radanliev@open.ac.uk}
}

\date{\today}

\begin{document}

\maketitle

\begin{abstract}
Open-weight machine learning models distributed via public repositories increasingly underpin critical artificial intelligence systems. Model cards serve as the primary public transparency mechanism communicating behavioral boundaries, biases, and safety evaluations. Here, we conduct a comprehensive empirical evaluation of safety disclosure across 809,703 model cards and 48,247 supply-chain lineage edges on the Hugging Face Hub, encompassing 2.47 billion cumulative downloads. Contrary to the assumption that widely adopted models carry higher documentation standards, we find that the download-weighted popularity gap is falsified on the registered estimator ($-2.9$ percentage points). On the unregistered row-level vector, the download-weighted safety rate (38.40\%) exceeds the unweighted rate (27.19\%) by 11.21 points, but the 95\% bootstrap interval ([$ -1.67$, $+24.68$] points) spans both confirmation and falsification boundaries. Extreme Value Theory reveals that model downloads follow an infinite-variance Pareto distribution (Hill tail index $\alpha = 0.98 \le 2$), rendering global weighted safety metrics hypersensitive to individual corporate releases: omitting the single most popular model shifts the global rate by 7.20 percentage points. Furthermore, derivative lineages exhibit severe documentation decay, with an 83.46\% license omission rate among model merges. These findings demonstrate that unadjusted download weighting yields fragile governance indicators, highlighting the need for tail-trimmed estimators and structural provenance tracking.
\end{abstract}

\section*{Introduction}
Public model hubs have become central distribution infrastructure for foundation artificial intelligence models. While commercial APIs incorporate centralized inference filters, open-weight deployments delegate risk assessment directly to developers. In this decentralised ecosystem, model cards provide the formal medium for documenting safety disclosures, intended uses, and ethical boundaries.

\section*{Results}
\subsection*{Empirical Census and Falsification of the Popularity Dividend}
Across 809,703 model cards, baseline safety disclosure stands at 27.19\%. On the registered binned estimator, the popularity gap yields $-2.9$ percentage points, falsifying the hypothesis that high-traffic models carry substantively higher documentation standards. The unregistered row-level estimate yields $+11.21$ points, with an indeterminate 95\% bootstrap interval of [$-1.67$, $+24.68$] points.

\subsection*{Extreme Value Theory and Super-Spreader Sensitivity}
Hill tail estimation confirms an infinite-variance Pareto regime ($\alpha = 0.98$). Influence function analysis demonstrates extreme leverage from individual corporate checkpoints: leaving out the top-downloaded model swings the global rate by 7.20 points, while trimming the top ten models compresses the observed gap to $+1.34$ points.

\subsection*{Documentation Decay in Derivative Supply Chains}
Tracing 48,247 lineage edges reveals that downstream fine-tunes and merges fail to maintain upstream safety disclosures. Model merges exhibit an 83.46\% license omission rate, creating systemic legal and technical documentation debt.

\section*{Discussion}
Our findings indicate that market mechanisms alone do not ensure reliable AI transparency on public registries. Policy frameworks such as the EU AI Act must account for heavy-tailed consumption dynamics and mandate verifiable, tail-trimmed accountability reporting.

\section*{Methods}
\subsection*{Data Extraction and Construct Validation}
We extracted card text and frontmatter from snapshot fef7adc9. Cards were categorized into four ordinal levels ($L_0$ Absent, $L_1$ Boilerplate, $L_2$ Qualitative, $L_3$ Quantitative). Human adjudication over stratified validation subsets confirmed high extractor reliability ($\kappa = 0.81$).

\subsection*{Extreme Value and Bootstrap Analysis}
Pareto tail indices were computed using the Hill estimator over the upper order statistics. Bootstrap confidence intervals were derived using 1,000 resamples across the empirical Lorenz curve.

\bibliographystyle{naturemag}
\bibliography{manuscript}

\end{document}
